\chapterimage{figs/chapterhead.png}
\chapter{Development Environment and Build}
\label{chap:development-environment-and-build}

This chapter describes the current build surface for GenESyS as it exists in
the repository now. It is based on the active root README, the current
\texttt{CMakePresets.json}, the local-agent runbook, and the focused sanitizer
path that already exists in \texttt{source/tests/}. The chapter does not
invent a release workflow that the repository does not currently expose.

\noindent\textbf{Objective.} Describe the environment, presets, target families,
and validation paths used by developers building GenESyS from source.

\noindent\textbf{Audience.} Contributors, maintainers, and anyone who needs to
configure, build, and check the project locally.

\noindent\textbf{Scope.} Cover the toolchain baseline, the current preset
matrix, separate build trees, debug and sanitizer posture, and the current
cleanup/incremental-build expectations.

\noindent\textbf{Status.} Confirmed by repository inspection and local command
execution on 2026-07-23. The chapter records the current build surface rather
than a hypothetical future release configuration.

The chapter follows this sequence:
Figure~\ref{fig:dev-build-pipeline},
Figure~\ref{fig:dev-preset-matrix},
Figure~\ref{fig:dev-target-families}, and
Figure~\ref{fig:dev-dependency-tree}.

\section{Baseline toolchain}
\label{sec:dev-baseline-toolchain}

The documented platform baseline is Ubuntu 24.04 with \texttt{CMake}
3.24 or newer, \texttt{Ninja}, \texttt{C++23} with compiler extensions
disabled, and Qt6 for the graphical applications. The current
\texttt{CMakePresets.json} encodes that policy directly by requiring
\texttt{CMAKE\_CXX\_STANDARD 23}, \texttt{CMAKE\_CXX\_STANDARD\_REQUIRED ON},
and \texttt{CMAKE\_CXX\_EXTENSIONS OFF}.

The local validation snapshot captured the following versions:
\begin{itemize}
\item \texttt{cmake} 3.28.3;
\item \texttt{ninja} 1.11.1;
\item \texttt{g++} 13.3.0;
\item Qt 6.4.2 via \texttt{qtpaths6 --qt-version}.
\end{itemize}

The repository root README still describes the intended baseline as
\texttt{CMake} 3.24 or newer, \texttt{Ninja}, \texttt{C++23}, Qt6, and the
configured Google Test fallback. That is the policy baseline; the versions
above are the observed local environment used for this chapter revision.

\subsection{Why the build surface is split}
\label{subsec:dev-why-build-surface-split}

GenESyS is not built as one monolith with a single universal preset. The
current preset matrix separates tests, shell, worker, GUI, and
model-specific targets into distinct build trees so that each workflow can be
configured and rebuilt independently.

This separation matters because different development tasks need different
dependencies and different validation scopes. A shell or model-specific build
does not need the same GUI surface as \texttt{gui-app}; a focused unit or
sanitizer check does not need the entire application set.

\begin{figure}[htbp]
\centering
% FIGURE-SPEC-BEGIN
% ID: fig:dev-build-pipeline
% Source of truth:
%   - ../README.md
%   - ../docs/ai_assistants/runbooks/LOCAL_AGENT.md
%   - ../CMakePresets.json
%   - ../source/tests/CMakeLists.txt
% Purpose:
%   Show the current configure -> build -> test flow for the main preset families.
% Required visual content:
%   - repository baseline and version checks
%   - configure step with a named preset
%   - build step with the matching build preset
%   - test step when the preset family has a test preset
% Accessibility description:
%   A pipeline diagram that explains how a developer moves from toolchain check to configure, build, and validation.
% Validation:
%   The command names and arrows must match the current presets and runbook guidance.
% FIGURE-SPEC-END
\figureplaceholder{figs/generated/dev-build-pipeline}
\caption{Current configure, build, and test pipeline.}
\label{fig:dev-build-pipeline}
\end{figure}

Figure~\ref{fig:dev-build-pipeline} is the canonical view of the local
developer workflow. It should make it obvious that configure, build, and test
are separate steps and that the test step only exists where the preset family
defines it.

\section{Current preset matrix}
\label{sec:dev-current-preset-matrix}

The current \texttt{cmake --list-presets=all} output is grouped into configure,
build, and test presets. The configure presets are:
\begin{itemize}
\item \texttt{tests-unit};
\item \texttt{tests-kernel-unit};
\item \texttt{genesys\_shell};
\item \texttt{terminal-smart};
\item \texttt{terminal-model-specific};
\item \texttt{genesys\_modelspecific\_app};
\item \texttt{terminal-smart-hold-search-remove};
\item \texttt{terminal-smart-ai-assistant};
\item \texttt{worker-app};
\item \texttt{genesys\_worker\_app};
\item \texttt{gui-app};
\item \texttt{gui-httpworker};
\item \texttt{gui-dataanalyser};
\item \texttt{gui-optimizer};
\item \texttt{gui-ai-assistant};
\item \texttt{tests-smoke}.
\end{itemize}

The build presets mirror the same names, with a special case for
\texttt{tests-kernel-unit-run} as the command-side alias for the kernel unit
test runner. The test presets are currently limited to
\texttt{tests-unit}, \texttt{tests-kernel-unit}, and \texttt{tests-smoke}.

All named configure presets currently inherit \texttt{debug-base}. That means
the current named surface is debug-oriented. There is no dedicated release
preset in the current inventory, so release builds must be introduced and
revalidated explicitly if they are needed later.

\begin{figure}[htbp]
\centering
% FIGURE-SPEC-BEGIN
% ID: fig:dev-preset-matrix
% Source of truth:
%   - ../CMakePresets.json
%   - ../README.md
% Purpose:
%   Present the current configure, build, and test preset families with their build trees.
% Required visual content:
%   - configure preset names
%   - matching build preset names
%   - matching test preset names
%   - the separate build directory for each family
% Accessibility description:
%   A matrix or table that lets the reader choose the correct preset without guessing.
% Validation:
%   The displayed preset names must match the current preset inventory.
% FIGURE-SPEC-END
\figureplaceholder{figs/generated/dev-preset-matrix}
\caption{Current preset matrix and build-tree separation.}
\label{fig:dev-preset-matrix}
\end{figure}

Figure~\ref{fig:dev-preset-matrix} should make the current split explicit:
tests use \path{build/tests-unit}, \path{build/tests-kernel-unit}, and
\path{build/tests-smoke}; the shell uses \path{build/genesys_shell}; the
worker uses \path{build/genesys_worker_app}; the GUI uses
\path{build/gui-app}; and the model-specific variants use their own build
trees under \path{build/}.

\subsection{Main build commands}
\label{subsec:dev-main-build-commands}

The repository root README and the local-agent runbook both use the same
command shape for configure and build. The current baseline examples are:
\begin{itemize}
\item \texttt{cmake --preset tests-unit};
\item \texttt{cmake --build --preset tests-unit --parallel "\$(nproc)"};
\item \texttt{ctest --preset tests-unit --output-on-failure};
\item \texttt{cmake --preset tests-kernel-unit};
\item \texttt{cmake --build --preset tests-kernel-unit --parallel "\$(nproc)"};
\item \texttt{ctest --preset tests-kernel-unit --output-on-failure};
\item \texttt{cmake --preset tests-smoke};
\item \texttt{cmake --build --preset tests-smoke --parallel "\$(nproc)"};
\item \texttt{ctest --preset tests-smoke --output-on-failure}.
\end{itemize}

Application builds follow the same pattern but target the application preset:
\begin{itemize}
\item \texttt{cmake --preset genesys\_shell};
\item \texttt{cmake --build --preset genesys\_shell --parallel "\$(nproc)"};
\item \texttt{cmake --preset genesys\_worker\_app};
\item \texttt{cmake --build --preset genesys\_worker\_app --parallel "\$(nproc)"};
\item \texttt{cmake --preset gui-app};
\item \texttt{cmake --build --preset gui-app --parallel "\$(nproc)"}.
\end{itemize}

The same pattern applies to the independent GUI and model-specific presets.
The build tree remains separate for each preset, so changing one family does
not overwrite another family unless the developer deliberately reuses the same
binary directory.

\section{Target families}
\label{sec:dev-target-families}

The current preset set maps onto a small number of real target families:
\begin{itemize}
\item regression tests: \texttt{genesys\_kernel\_unit\_tests},
\texttt{genesys\_kernel\_unit\_tests\_run}, and \texttt{genesys\_smoke\_tests};
\item shell: \texttt{genesys\_shell};
\item model-specific terminal application: \texttt{genesys\_terminal\_application};
\item dedicated model-specific app: \texttt{genesys\_modelspecific\_app};
\item worker: \texttt{genesys\_worker\_app};
\item main GUI: \texttt{genesys\_gui};
\item independent GUI frontends:
\begin{itemize}
\item \texttt{genesys\_httpworker\_gui\_application};
\item \texttt{genesys\_dataanalyser\_gui\_application};
\item \texttt{genesys\_optimizer\_gui\_application};
\item \texttt{genesys\_ai\_assistant\_gui\_application}.
\end{itemize}
\end{itemize}

The shell and model-specific presets also encode source selection through
cache variables such as \texttt{GENESYS\_TERMINAL\_EXAMPLE} and
\texttt{GENESYS\_MODELSPECIFIC\_APP}. That means the source tree selection is
part of the build contract, not just an ad hoc command-line argument.

\begin{figure}[htbp]
\centering
% FIGURE-SPEC-BEGIN
% ID: fig:dev-target-families
% Source of truth:
%   - ../CMakePresets.json
%   - ../README.md
% Purpose:
%   Show which logical application family produces which build target.
% Required visual content:
%   - regression test family
%   - shell family
%   - worker family
%   - main GUI family
%   - independent GUI family
% Accessibility description:
%   A family map that lets the reader see which targets belong to each executable area.
% Validation:
%   Each family label must match a real target name in the current tree.
% FIGURE-SPEC-END
\figureplaceholder{figs/generated/dev-target-families}
\caption{Current target families and their executable boundaries.}
\label{fig:dev-target-families}
\end{figure}

Figure~\ref{fig:dev-target-families} should help developers avoid the common
mistake of building the wrong family when the task only needs a single target
or one independent GUI.

\section{Dependency and decision tree}
\label{sec:dev-dependency-tree}

The current repository does not expose every possible build choice as a named
preset. The conservative decision tree is simple:
\begin{enumerate}
\item If the task is about ordinary regression, use one of the test presets.
\item If the task is about the interactive shell, use \texttt{genesys\_shell}.
\item If the task is about the worker runtime, use \texttt{genesys\_worker\_app}.
\item If the task is about the main GUI, use \texttt{gui-app}.
\item If the task is about an independent GUI, use the matching GUI preset.
\item If the task is about model-specific generation, use the matching
terminal/model-specific preset.
\item If the task needs a sanitizer, use the focused sanitizer-enabled test
executable that already exists, rather than assuming a whole-repository
sanitizer preset.
\end{enumerate}

The current sanitizer path is intentionally focused. In
\texttt{source/tests/CMakeLists.txt}, the
\texttt{genesys\_test\_simulator\_plugin\_completion\_lifetime} target adds
\texttt{-fsanitize=address} and \texttt{-fno-omit-frame-pointer} under GNU or
Clang, links with the matching sanitizer option, and runs with
\texttt{ASAN\_OPTIONS=detect\_leaks=1:halt\_on\_error=1} and
\texttt{LSAN\_OPTIONS=exitcode=23}. That proves a focused ownership/lifetime
sanity path exists. It does not prove repository-wide leak freedom.

The chapter's dependency diagram makes that choice explicit.

\begin{figure}[htbp]
\centering
% FIGURE-SPEC-BEGIN
% ID: fig:dev-dependency-tree
% Source of truth:
%   - ../README.md
%   - ../CMakePresets.json
%   - ../source/tests/CMakeLists.txt
% Purpose:
%   Show how a developer chooses between test, shell, worker, GUI, model-specific, and sanitizer builds.
% Required visual content:
%   - root toolchain check
%   - branch for tests
%   - branch for shell/worker/GUI/model-specific apps
%   - branch for focused sanitizer execution
% Accessibility description:
%   A decision tree that points the reader to the right preset and validation path.
% Validation:
%   The branch names and target names must match the current repository surface.
% FIGURE-SPEC-END
\figureplaceholder{figs/generated/dev-dependency-tree}
\caption{Current build dependency decision tree.}
\label{fig:dev-dependency-tree}
\end{figure}

Figure~\ref{fig:dev-dependency-tree} is the practical guide for selecting the
right preset. It also makes the sanitizer boundary explicit: there is a real
focused sanitizer path, but not a universal sanitizer preset.

\section{Warnings, cleanup, and incremental build}
\label{sec:dev-warnings-cleanup-incremental}

The current workflow assumes incremental rebuilds inside each preset-specific
build tree. That is the default \texttt{cmake --build} behavior: only changed
targets are rebuilt unless the developer asks for a fresh configuration.

Use a separate build tree when changing incompatible options. If the current
cache becomes stale, rebuild from the affected \path{build/<preset>} tree
instead of mixing unrelated preset state.

Warnings should be treated conservatively:
\begin{itemize}
\item configuration warnings should be resolved or explicitly understood;
\item build warnings should be checked against the target being changed;
\item a successful build does not by itself validate runtime behavior;
\item a successful startup does not by itself validate the workflow or the
science behind it.
\end{itemize}

The current preset inventory is debug-oriented, so a release-oriented build
requires a deliberate new configuration and a fresh validation pass. Likewise,
sanitizers are currently a focused diagnostic path rather than a repository
mode that can be assumed everywhere.

\section{Validation snapshot}
\label{sec:dev-validation-snapshot}

The local documentation pass for this chapter used the following commands:
\begin{itemize}
\item \texttt{cmake --list-presets=all};
\item \texttt{cmake --version};
\item \texttt{ninja --version};
\item \texttt{g++ --version};
\item \texttt{qtpaths6 --qt-version};
\item \texttt{source/tests/CMakeLists.txt} inspection for the focused sanitizer target.
\end{itemize}

The observed versions were \texttt{cmake} 3.28.3, \texttt{ninja} 1.11.1,
\texttt{g++} 13.3.0, and Qt 6.4.2. Those values are the current evidence
snapshot for this chapter revision and should be refreshed when the toolchain
changes.
